\section{Background}
\label{sec:background}

\subsection{Structured state memory}

\SkillState{}~\citep{skillstate2026} maintains a persistent structured state
$\Sigma_t$ and assembles each prompt as
\[
  \text{prompt}_t \;=\; [\,\text{task spec}\,] \;+\; [\,\Sigma_t\,] \;+\;
  [\,o_t\,],
\]
where $o_t$ is the latest observation and the interaction history is thrown
away. In one generation the model emits its reasoning, a state patch and an
action, $(R_t, \Delta\Sigma_t, a_t)$, as a single JSON object with exactly
the two keys \texttt{state\_patch} and \texttt{action}\origin. The runtime
checks each patch against a schema and rejects invalid ones with a
rollback-retry cycle. Schemas are written once per domain; across 100
InterCode CTF instances a single five-field schema is reused\origin.

The reported evidence is a horizon sweep on a procedural generator: the
prompt stays near \num{1800} characters from $T{=}10$ to $T{=}200$ while the
baselines grow to \num{48000}, and quality is preserved or improved\origin.
The paper's own error taxonomy for open-weight models attributes
\SI{68}{\percent} of failures to premature state overwrite or deletion\origin,
a figure we return to in \S\ref{sec:composition}, because it is an argument
for moving part of the state out of the model's hands.

\subsection{Deterministic execution state}

\Ledger{}~\citep{ledger2026} sits in the opposite design corner: a record
maintained by the runtime and derived mechanically from completed
interactions: observation records, per-file and global modification
counters, normalised command records with their outcomes. It adds no model
calls and does not modify the agent. Its own ablation separates its two
mechanisms cleanly: \emph{inform} (rendering the record at the end of the
input, where it cannot disturb an append-only prefix) produces most of the
efficiency gain, while \emph{govern} (intercepting a proposed command and
reusing an earlier result) produces most of the resolution gain\origin.

The two systems are complementary rather than competing: \SkillState{} keeps
semantic state in the model, \Ledger{} keeps objective state in the runtime.
Our composition study (\S\ref{sec:composition}) varies exactly this boundary,
and we deliberately exclude \emph{govern} from it: govern changes what the
agent does, so including it would mix a large effect into the measurement of
a smaller one.

\subsection{The assumptions under test}

Stated as checkable propositions rather than criticisms:

\begin{description}[style=nextline]
\item[A1 (schema availability)] An adequate domain schema exists before
  execution. The original gives no schema-design method and reports a
  verbatim field list for only one of its four domains\origin.
\item[A2 (positive from turn one)] Maintaining state is at worst neutral
  from the first turn. The shortest horizon reported is $T{=}10$, already
  longer than a large share of real coding tasks.
\item[A3 (sufficient statistic)] Nothing discarded is ever needed again.
  The original names this as a limitation\origin{} but does not measure it.
  This is the assumption this paper falsifies.
\item[M (measurement currency)] Cumulative token count is a valid proxy for
  cost. A full-text search of the original returns zero occurrences of
  cache, billing or monetary terminology\origin.
\end{description}

A1 and A3 are properties of the method. M is a property of how the method
was evaluated, and it is the one that hides the other two: a metric that
counts logical tokens can represent neither the cache discount that rewards
an append-only prefix nor the tool calls spent re-acquiring what the prompt
no longer carries.

\subsection{The design axis this paper uses}

Instead of comparing named architectures, we hold the mechanism fixed and
vary one integer: the number $\tailk$ of most recent verbatim exchanges kept
next to the state block.
\[
  \text{prompt}_t \;=\; [\,\text{task spec}\,] \;+\;
  [\,\text{last } \tailk \text{ exchanges}\,] \;+\; [\,\Sigma_t\,]
\]
At $\tailk{=}0$ this is exactly \SkillState{} as published. As $\tailk$
grows, the configuration approaches a transcript agent. The control is not a
point on this axis (it carries no state block at all) and is always reported
separately.

This buys three things a three-arm comparison cannot buy. It removes the
arbitrary constant that an intermediate ``hybrid'' arm would introduce; it
changes the question from \emph{which architecture} to \emph{where on the
axis the optimum lies}; and it makes the shape of the dependence visible,
which is what ultimately localises the damage (\S\ref{sec:capability}).
